% Section 1.4: 研究目标与主要贡献

本文以“面向算法代码生成的大语言模型指令微调与推理增强”为主题，围绕1.3所归纳的核心问题，设定如下总体研究目标：

\begin{enumerate}
  \item 构建算法代码生成任务下训练–评测一体化的标准化实验体系，确保不同方法在数据、提示与评测维度具备严格可比性；
  \item 在有限算力约束下，利用参数高效微调技术，实现基座模型在算法题上可执行正确性与生成稳定性的量化提升；
  \item 引入结构化推理引导机制，系统对比零样本思维链与少样本思维链的效果差异，揭示其在不同复杂度问题上的边际贡献与适用条件；
  \item 对模型失效模式开展归因分析，形成可复现的实证结论，为算法代码生成系统的优化方向提供依据。
\end{enumerate}

围绕上述目标，本文的主要工作与贡献可以概括为以下四点：

\paragraph{（1）标准化、可复现的端到端实验框架}
本文构建了覆盖数据处理、提示构造、模型推理、可执行性评测的完整实验流程，实现训练与评测环节的统一规范。在统一协议下对基座模型、高效微调、零样本推理、少样本推理等方案开展对照，为方法对比提供了公平、可靠的实验基础。

\paragraph{（2）参数高效微调在算法代码生成中的有效落地}
面向算法题任务特点，设计适配的微调样本构造、训练目标与参数配置策略，在不显著增加推理开销的前提下，显著提升代码执行通过率与生成稳定性。实证结果验证了参数高效微调在算法代码生成场景下的有效性与工程实用性。

\paragraph{（3）双路径思维链推理的系统对照与机理分析}
在微调基础上，构建零样本思维链与少样本思维链两种结构化推理范式，在统一评测集上开展严格对照。通过多维度分析，揭示不同推理引导方式在简单 / 复杂题目上的边际效应，明确少样本示例的增益边界与组织优势，形成具有参考价值的工程化结论。

\paragraph{（4）失效模式归因与可复用实验支撑}
从逻辑正确性、代码结构、边界处理等角度对典型错误进行分类归因，识别出高频失效模式。整套实验流程、提示模板与自动化评测工具具备良好的可复现性与可迁移性，可为后续相关研究与系统迭代提供直接支撑。